% Part: propositional-logic
% Chapter: syntax-and-semantics
% Section: introduction

\documentclass[../../../include/open-logic-section]{subfiles}

\begin{document}

\olfileid{pl}{syn}{int}

\olsection{تعارف}

قضیاتی منطق اوہناں !!{formula}s نال واسطہ رکھدی اے جیہڑے
!!{propositional variable}s توں قضیاتی رابطیاں
$\lnot$، $\land$، $\lor$، $\lif$ تے $\liff$ نال بنائے جاندے نیں۔ سہج
فہم طور تے، !!a{propositional variable}~$p$ اک جملے یا قضیے دی
نمائندگی کردا اے جیہڑا سچ یا جھوٹ ہوندا اے۔ جدوں وی !!a{formula} وچ
!!{propositional variable} دی ``صداقتی قدر'' متعین ہو جاوے، تاں
قضیاتی رابطیاں نال اوہناں توں بنے !!{formula}s وچوں ہر اک دی صداقتی قدر وی متعین
ہو جاندی اے۔ اسیں آکھدے آں کہ قضیاتی منطق
\emph{صداقتی-فنکشنی} اے، کیوں جے ایہدی معنیات صداقتی قدراں دے فنکشناں
راہیں دِتی جاندی اے۔ خاص طور تے، قضیاتی منطق وچ اسیں صدق تے کذب دی
ہور کوئی تعیین زیرِ غور نہیں لیاؤندے، مثلاً ایہ کہ کوئی گل محض اتفاقی
طور تے سچ ہون دی بجائے لازماً سچ اے، یا ایہ کہ کوئی گل سچ معلوم اے، یا
ایہ کہ کوئی گل ہُن سچ اے نہ کہ پہلے سچ سی یا اگوں سچ ہووے گی۔ اسیں صرف
دو صداقتی قدراں، سچ ($\True$) تے جھوٹ ($\False$)، اُتے غور کردے آں، تے
ایس لئی ایہ امکان بحث توں باہر رکھدے آں کہ کوئی بیان نہ سچ ہووے نہ
جھوٹ، یا صرف ادھا سچ ہووے۔ اسیں صرف اوہناں رابطیاں اُتے وی دھیان
دیندے آں جنہاں توں بنے !!a{formula} دی صداقتی قدر اوہدے حصیاں دیاں
صداقتی قدراں نال پوری طرح متعین ہوندی اے (نہ کہ، مثلاً، اوہدے معنی
نال)۔ خاص طور تے، انگریزی شرطیہ جملیاں دی صداقتی قدر ایس معنی وچ
صداقتی-فنکشنی اے یا نہیں، ایہ اختلافی گل اے۔ مادی شرطیہ~$\lif$
صداقتی-فنکشنی اے؛ ہور منطقاں اوہناں شرطیاں نال واسطہ رکھدیاں نیں
جیہڑیاں صداقتی-فنکشنی نہیں ہُندیاں۔

صداقتی-فنکشنی قضیاتی منطق دا نظریہ تے مابعد نظریہ تیار کرن لئی، سانوں
پہلاں ایہدیاں عبارتاں دی نحو تے معنیات دی تعریف کرنی پئے گی۔ اسیں
رابطیاں نال !!{propositional variable}s توں !!{formula}s بناؤن دا اک
طریقہ بیان کراں گے۔ ہور تعریفاں وی ممکن نیں۔ ہور نظام وکھریاں علامات
چُنن گے، رابطیاں دے وکھرے مجموعیاں نوں ابتدائی منّن گے، تے قوسین نوں
وکھرے ڈھنگ نال ورتن گے (یا نام نہاد پولش علامت نگاری وانگ بالکل نہیں
ورتن گے)۔ پر سارے طریقیاں وچ سانجھی گل ایہ اے کہ تشکیل دے قاعدے
!!{formula}s دے سیٹ دی
\emph{استقرائی طور تے} تعریف کردے نیں۔ جے ایہ کم ٹھیک طرح کیتا جاوے،
تاں ہر عبارت بنیادی طور تے تشکیل دے قاعدیاں مطابق صرف اکّو طریقے نال
بن سکدی اے۔ فارمولیاں نوں
\emph{اکّو طرح نال پڑھن جوگ} بناؤن والی استقرائی تعریف دا مطلب اے کہ
اسیں ایسّے طریقے، یعنی استقرائی تعریف، نال ایہناں عبارتاں دے معنی وی
دے سکدے آں۔

عبارتاں نوں معنی دینا معنیات دا دائرہ اے۔ قضیاتی منطق دی معنیات وچ
مرکزی تصور !!a{valuation} وچ پورا اترن دا اے۔
!!^a{valuation}~$\pAssign{v}$، !!{propositional variable}s نوں صداقتی
قدراں $\True$، $\False$ مقرر کردی اے۔ ہر !!{valuation} اک صداقتی قدر
$\pValue{v}(!A)$، کسی وی !!{formula}~$!A$ لئی، متعین کردی اے۔
!!^a{formula}، !!a{valuation}~$\pAssign{v}$ وچ اُدوں تے صرف اُدوں پورا
اتردا اے جدوں $\pValue{v}(!A) = \True$ ہووے—اسیں ایہنوں
$\pSat{v}{!A}$ لکھدے آں۔ ایس تعلق دی تعریف $!A$ دی ساخت اُتے استقرا
نال وی کیتی جا سکدی اے؛ منطقی رابطیاں دے صداقتی فنکشن ورت کے، مثلاً،
$!A \land !B$ دے پورا اترن نوں $!A$ تے~$!B$ دے پورا اترن (یا نہ
اترن) دے لحاظ نال تعریف کیتا جا سکدا اے۔

جملیاں لئی پورا اترن دے تعلق $\pSat{v}{!A}$ دی بنیاد اُتے اسیں فیر
ہمیشہ صادق ہون، منطقی نتیجے تے پورا ہون جوگتا دے بنیادی معنوی تصوراں
دی تعریف کر سکدے آں۔ !!^a{formula} اک ہمیشہ صادق فارمولا،
$\Entails !A$، اے جے ہر !!{valuation} اوہنوں پورا کردی اے، یعنی
$\pValue{v}(!A) = \True$ ہر~$\pAssign{v}$ لئی۔ اوہی فارمولا
!!{formula}s دے سیٹ دا منطقی نتیجہ ہوندا اے، $\Gamma \Entails !A$، جے ہر اوہ
!!{valuation} جیہڑی~$\Gamma$ دے سارے !!{formula}s نوں پورا کردی اے،
اوہ~$!A$ نوں وی پورا کردی اے۔ تے !!{formula}s دا سیٹ پورا ہون جوگ اے
جے کوئی !!{valuation} اوس وچلے سارے !!{formula}s نوں اکّو ویلے پورا
کردی اے۔ کیوں جے !!{formula}s دی تعریف استقرائی طور تے کیتی گئی اے،
تے پورا اترن دی تعریف اگوں !!{formula}s دی ساخت اُتے استقرا نال کیتی
گئی اے، ایس لئی اسیں اپنی معنیات دیاں خاصیتاں ثابت کرن تے تعریف کیتے
معنوی تصوراں نوں آپس وچ جوڑن لئی استقرا ورت سکدے آں۔


\end{document}
